% !TeX root = ../main.tex

% ==================================================
% 2.3. 클래스와 인스턴스 — Person과 Person()은 어떻게 다른가?
% ==================================================

\subsection{클래스와 인스턴스 — \zsh{Person}과 \zsh{Person()}은 어떻게 다른가?}

2.2에서는 Python의 \zsh{class}문이 단순한 선언이 아니라, 클래스 본문을 실행하고 클래스 객체를 생성한 뒤 이름에 바인딩하는 문장이라는 사실을 학습했다.

이번 2.3에서는 생성된 클래스 객체를 호출하면 어떤 일이 발생하는지 살펴본다.

특히 다음 세 표현식의 차이를 정확하게 이해하는 것이 목표다.

\begin{dhkcode}
Person
Person()
p = Person()
\end{dhkcode}

겉으로 보기에는 \zsh{()} 의 유무와 대입문 정도의 차이지만,

Python의 객체 모델에서는 서로 다른 동작을 의미한다.

이번 학습에서는 다음 질문에 답할 것이다.

\begin{dhksmry}
1. 클래스와 인스턴스는 정확히 무엇이 다른가?

2. \zsh{Person}과 \zsh{Person()}은 각각 무엇을 의미하는가?

3. 인스턴스가 생성될 때 클래스 객체도 복사되는가?

4. 하나의 클래스로 여러 인스턴스를 생성하면 어떤 관계가 형성되는가?

5. 인스턴스는 자신의 클래스를 어떻게 알 수 있는가?

6. \zsh{is}, \zsh{==}, \zsh{type()}, \zsh{isinstance()}는 어떤 차이가 있는가?

7. 이러한 구조가 Pytorch 모델에서는 어떻게 사용되는가?
\end{dhksmry}

\subsubsection{개념 정의 - 클래스와 인스턴스}

먼저 가장 기본적인 코드부터 살펴보자.

\begin{dhkcode}
class Person:
    pass
\end{dhkcode}

이 코드를 실행하면 Python은 \zsh{Person}이라는 클래스 객체를 생성한다.

여기서 아직 \zsh{Person}의 인스턴스를 생성한 것은 아니다.

이후 다음 코드를 실행한다고 가정하자.

\begin{dhkcode}
p = Person()
\end{dhkcode}

이때 일반적인 클래스에서는 새로운 인스턴스 객체가 생성된다.

따라서 두 객체를 구분해야 한다.

\begin{dhknote}
\begin{dhksmry}[centered]
클래스 객체(Class Object)
\end{dhksmry}
\end{dhknote}

클래스 객체는 인스턴스의 생성과 동작을 정의하는 객체다.

\begin{dhkcode}
Person
\end{dhkcode}

\zsh{Person}이라는 이름은 클래스 객체를 참조한다.

\begin{dhknote}
\begin{dhksmry}[centered]
인스턴스 객체(Instance Object)
\end{dhksmry}
\end{dhknote}

인스턴스는 클래스 객체를 기반으로 생성된 개별 객체다.

\begin{dhkcode}
Person()
\end{dhkcode}

이 표현식은 클래스 객체를 호출한다.

일반적인 경우 새로운 인스턴스를 생성하고 반환한다.

\begin{dhknote}
\begin{dhksmry}[centered]
인스턴스를 참조하는 이름
\end{dhksmry}
\end{dhknote}

\begin{dhkcode}
p = Person()
\end{dhkcode}

오른쪽 표현식의 평가 결과를 왼쪽 이름 \zsh{p}에 바인딩한다.

따라서 전체 관계는 다음과 같다.

\begin{dhkdiagram}[in=6]
node per (1,1) :: Person
node c (1,3) :: 클래스 객체
node i (3,3) :: 인스턴스 객체
node p (5,1) :: p

edge per -> c
edge c -> i :: Person()호출
edge p -> i :: HV
\end{dhkdiagram}

여기서 중요한 사실은 클래스 객체와 인스턴스 객체가 별도로 존재한다는 것이다.

클래스는 인스턴스를 생성하기 위한 정의이지만, Python에서는 그 정의 자체도 하나의 객체다.

\begin{dhktable}
|구분|	클래스 객체|	인스턴스 객체|
|:---|:---|:---|
|예시|	Person|	Person()의 반환값|
|일반적인 타입|	type|	Person|
|생성 시점|	class 정의문 실행 시|	클래스 호출 시|
|주요 역할|	인스턴스의 구조와 동작 정의|	개별 객체의 상태와 동작 표현|
|여러 개 생성 가능 여부|	클래스 정의를 다시 실행하면 새 클래스 객체 생성 가능|	같은 클래스에서 여러 인스턴스 생성 가능|
\end{dhktable}

여기서 인스턴스가 반드시 개별적인 상태를 가져야 하는 것은 아니다.

\begin{dhkcode}
class Person:
    pass
\end{dhkcode}

이처럼 속성이 없는 클래스의 인스턴스도 정상적인 인스턴스 객체다.

\subsubsection{\zsh{Person}과 \zsh{Person()}의 차이}

이제 두 표현식을 비교해 보자.

\begin{dhkcode}
Person

Person()
\end{dhkcode}

이전에 함수 객체를 공부하면서 다음 두 표현식의 차이를 학습했다.

\begin{dhkcode}
add

add()
\end{dhkcode}

\zsh{add}는 함수 객체를 참조하고, \zsh{add()}는 함수 객체를 호출하는 표현식이다.

클래스도 이와 비슷하다.

\begin{dhknote}
\begin{dhksmry}[centered]
Person - 클래스 객체 참조
\end{dhksmry}
\end{dhknote}

\begin{dhkcode}
class Person:
    pass

print(Person)
\end{dhkcode}

출력:

\begin{dhkzsh}
<class '__main__.Person'>
\end{dhkzsh}

여기서는 클래스 객체를 참조할 뿐, 새로운 인스턴스를 생성하지 않는다.

\begin{dhknote}
\begin{dhksmry}[centered]
Person() - 클래스 객체 호출
\end{dhksmry}
\end{dhknote}

\begin{dhkcode}
p = Person()
\end{dhkcode}

\zsh{()}는 앞에 있는 호출 가능한 객체를 호출하는 문법이다.

따라서 Python은 \zsh{Person}이 참조하는 클래스 객체를 호출한다.

일반적인 클래스 객체 호출은 인스턴스 생성과 초기화 과정으로 이어진다.

\begin{dhknote}
\begin{dhksmry}[centered]
p = Person() - 반환값 바인딩
\end{dhksmry}
\end{dhknote}

\begin{dhkcode}
p = Person()
\end{dhkcode}

이 대입문을 개념적으로 나누면 다음과 같다.

\begin{dhksmry}[centered]
1. Person 이름 평가

\da

2. 클래스 객체 획득

\da

3. 클래스 객체 호출

\da

4. 인스턴스 생성 및 초기화

\da

5. 호출 결과 반환

\da

6. p 이름에 반환된 객체 바인딩
\end{dhksmry}

중요한 것은 \zsh{p}가 객체를 생성하는 것이 아니라는 점이다.

객체 생성은 오른쪽의 \zsh{Person()} 호출 과정에서 이루어진다.

왼쪽의 \zsh{p}는 그 결과를 참조하는 이름일 뿐이다.

이전에 학습한 이름과 객체의 바인딩 개념이 그대로 적용되는 것이다.

\subsubsection{하나의 클래스에서 여러 인스턴스 생성하기}

다음 코드를 살펴보자.

\begin{dhkcode}
class Person:
    pass

p1 = Person()
p2 = Person()
p3 = Person()
\end{dhkcode}

여기서 \zsh{Person} 클래스 객체는 하나다.

하지만 \zsh{Person()}을 세 번 호출했으므로 일반적인 생성 과정에서는 서로 다른 인스턴스 세 개가 만들어진다.

\begin{dhkdiagram}[in=6]
node per (1,2) :: Person
node c (3,2) :: 클래스 객체
node pi1 (4,1) :: Person()
node pi2 (4,2) :: Person()
node pi3 (4,3) :: Person()
node i1 (6,1) :: 인스턴스1
node i2 (6,2) :: 인스턴스2
node i3 (6,3) :: 인스턴스3
node p1 (8,1) :: p1
node p2 (8,2) :: p2
node p3 (8,3) :: p3

edge per -> c
split c -> pi1, pi2, pi3
edge pi1 -> i1
edge pi2 -> i2
edge pi3 -> i3

edge p1 -> i1
edge p2 -> i2
edge p3 -> i3
\end{dhkdiagram}

실제로 확인해 보자.

\begin{dhkcode}
print(p1 is p2)
print(p2 is p3)
print(p1 is p3)
\end{dhkcode}

출력:

\begin{dhkzsh}
False
False
False
\end{dhkzsh}

\zsh{is}는 두 표현식이 동일한 객체를 참조하는지 검사한다.

각각 새로운 인스턴스를 생성했기 때문에 결과가 모두 \zsh{False}다.

그런데 다음 코드는 어떨까?

\begin{dhkcode}
p1 = Person()
p2 = p1

print(p1 is p2)
\end{dhkcode}

출력:

\begin{dhkzsh}
True
\end{dhkzsh}

이번에는 \zsh{Person()}을 한 번만 호출했다.

\zsh{p2 = p1}은 새로운 인스턴스를 생성하는 코드가 아니다.

기존 인스턴스에 새로운 이름을 바인딩한다.

\begin{dhkdiagram}[in=6]
node p1 (1,1) :: p1
node i (2,2) :: 인스턴스 객체
node p2 (3,1) :: p2

edge p1 -> i :: HV
edge p2 -> i :: HV
\end{dhkdiagram}

이 구조는 이전에 살펴본 이름과 객체의 관계와 동일하다.

이름을 새로 만든다고 객체가 새로 생성되는 것은 아니다.

\subsubsection{인스턴스가 생성되면 클래스도 복사되는가?}

여기서 중요한 질문이 하나 생긴다.

\begin{dhkcode}
p1 = Person()
p2 = Person()
\end{dhkcode}

이렇게 인스턴스를 두 개 생성하면 \zsh{Person} 클래스 객체도 두 개로 복사될까?

아니다.

일반적인 인스턴스 생성에서는 클래스 객체가 복사되지 않는다.

두 인스턴스는 동일한 클래스 객체를 자신의 타입으로 가진다.

다음 코드로 확인할 수 있다.

\begin{dhkcode}
class Person:
    pass

p1 = Person()
p2 = Person()

print(type(p1) is Person)
print(type(p2) is Person)
print(type(p1) is type(p2))
\end{dhkcode}

출력:
\begin{dhkzsh}
True
True
True
\end{dhkzsh}

관계를 표현하면 다음과 같다.

\begin{dhkdiagram}[in=6]
node per (1,2) :: Person
node c (3,2) :: 클래스 객체
node t1 (5,1) :: type
node t2 (5,3) :: type
node i1 (7,1) :: 인스턴스1
node i2 (7,3) :: 인스턴스2
node p1 (9,1) :: p1
node p2 (9,3) :: p2

edge per -> c
edge t1 -> c :: VH
edge t2 -> c :: VH
edge i1 -> t1
edge i2 -> t2
edge p1 -> i1
edge p2 -> i2
\end{dhkdiagram}

두 인스턴스는 각각 독립적인 객체이지만 동일한 클래스 객체를 타입으로 참조한다.

이 구조가 중요한 이유는 메서드와 클래스 속성 때문이다.

예를 들어 다음 클래스를 보자.

\begin{dhkcode}
class Person:

    def introduce(self):
        print("Hello")
\end{dhkcode}

인스턴스를 여러 개 생성해도 일반적으로 \zsh{introduce} 함수 객체가 인스턴스마다 복사되는 것은 아니다.

\begin{dhkcode}
p1 = Person()
p2 = Person()
\end{dhkcode}

\zsh{introduce} 함수 객체는 클래스의 네임스페이스에 존재한다.

\begin{dhkcode}
print(Person.__dict__["introduce"])
\end{dhkcode}

각 인스턴스에서 다음처럼 접근할 수 있다.

\begin{dhkcode}
p1.introduce()
p2.introduce()
\end{dhkcode}

Python은 인스턴스를 통해 메서드에 접근할 때 클래스의 함수 객체를 찾아 해당 인스턴스에 바인딩된 메서드로 제공한다.

이 과정은 2.5에서 자세히 학습한다.

현재 단계에서는 다음 사실을 기억하면 충분하다.

\begin{dhksmry}[centered]
인스턴스가 여러 개 존재하더라도 

클래스에 정의된 메서드가 인스턴스마다 복사되는 것은 아니다.
\end{dhksmry}

\subsubsection{인스턴스는 자신의 클래스를 어떻게 알 수 있는가?}

앞에서 다음 코드를 사용했다.

\begin{dhkcode}
type(p1)
\end{dhkcode}

그렇다면 Python은 \zsh{p1}이 \zsh{Person}의 인스턴스라는 사실을 어떻게 알까?

Python 객체는 자신의 타입에 대한 정보를 가진다.

일반적인 Python 객체에서는 \zsh{__class__} 속성을 통해 해당 정보를 확인할 수 있다.

\begin{dhkcode}
class Person:
    pass

p = Person()

print(p.__class__)
\end{dhkcode}

출력:

\begin{dhkzsh}
<class '__main__.Person'>
\end{dhkzsh}

다음 코드도 확인해 보자.

\begin{dhkcode}
print(p.__class__ is Person)
\end{dhkcode}

출력:

\begin{dhkzsh}
True
\end{dhkzsh}

따라서 다음 두 표현식은 이 예제에서 같은 클래스 객체를 반환한다.

\begin{dhkcode}
type(p)
p.__class__
\end{dhkcode}

다만 엄밀하게는 두 표현식의 동작이 항상 완전히 동일한 것은 아니다.

\zsh{type(p)}는 객체의 실제 타입을 확인하는 반면,

\zsh{p.__class__}는 속성 접근을 통해 얻는 값이므로 사용자 정의 클래스의 특수한 속성 처리에 영향을 받을 수 있다.

일반적인 Python 클래스를 공부하는 현재 단계에서는 둘 다 인스턴스와 클래스의 관계를 확인하는 방법으로 이해해도 충분하다.

\begin{dhknote}
\begin{dhksmry}[centered]
타입과 인스턴스 관계
\end{dhksmry}
\end{dhknote}

다음 코드를 다시 보자.

\begin{dhkcode}
class Person:
    pass

p = Person()
\end{dhkcode}

각 표현식의 타입을 확인하면 다음과 같다.

\begin{dhkcode}
print(type(p))
pirnt(type(Person))
\end{dhkcode}

출력:

\begin{dhkzsh}
<class '__main__.Person'>
<class 'type'>
\end{dhkzsh}

즉,

\begin{dhkzsh}
type(p) = Person
type(Person) = type
\end{dhkzsh}


이 관계는 2.2에서 배운 내용과 연결된다.

\begin{dhkdiagram}[in=6]
node t (1, 1) :: type
node per (3, 1) :: Person
node p (5, 1) :: p

edge t -> per :: 클래스 객체의 타입
edge per -> p :: 인스턴스 객체의 타입
\end{dhkdiagram}

여기서 다시 강조할 점은 이 관계가 상속 관계가 아니라 타입과 인스턴스의 관계라는 것이다.

상속은 이후 2.8에서 별도로 학습한다.

\subsubsection{객체의 동일성과 동등성 - \zsh{is}와 \zsh{==}}

클래스와 인스턴스를 이해할 때 반드시 구분해야 하는 개념이 있다.

바로 객체의 동일성(Identity)과 동등성(Equality)이다.

이전에 객체와 바인딩을 공부하면서 살펴봤던 개념이기도 하다.

\begin{dhknote}
\begin{dhksmry}[centered]
is - 동일한 객체인가?
\end{dhksmry}
\end{dhknote}

\begin{dhkcode}
class Person:
    pass

p1 = Person()
p2 = Person()
p3 = p1

print(p1 is p2)
print(p1 is p3)
\end{dhkcode}

출력:

\begin{dhkzsh}
False
True
\end{dhkzsh}

\zsh{is}는 두 이름이 같은 객체를 참조하는지 확인한다.

\begin{dhknote}
\begin{dhksmry}[centered]
== - 두 객체가 동등한가?
\end{dhksmry}
\end{dhknote}

\begin{dhkcode}
print(p1 == p2)
print(p1 == p3)
\end{dhkcode}

출력:

\begin{dhkzsh}
False
True
\end{dhkzsh}

이 예제에서는 \zsh{Person} 클래스에 동등성 비교를 위한 별도의 메서드를 정의하지 않았다.

따라서 기본적인 객체 비교 동작에 따라 서로 다른 인스턴스는 동등하지 않은 것으로 평가된다.

하지만 \zsh{==}의 동작은 클래스에서 \zsh{__eq__()}를 정의하여 변경할 수 있다.

다음 예제를 보자.

\begin{dhkcode}
class Person:
    def __init__(self, name):
        self.name = name

    def __eq__(self, other):
        if not isinstance(other, Person):
            return NotImplemented
        
        return self.name == other.name
\end{dhkcode}

이제 다음 코드를 실행해 보자.

\begin{dhkcode}
p1 = Person("Kim")
p2 = Person("Kim")

print(p1 is p2)
print(p1 == p2)
\end{dhkcode}

출력:

\begin{dhkzsh}
False
True
\end{dhkzsh}

왜 이런 결과가 발생할까?

\zsh{p1}과 \zsh{p2}는 서로 다른 인스턴스 객체다.

따라서 \zsh{is}는 \zsh{False}다.

하지만 \zsh{__eq__()}에서는 두 인스턴스의 \zsh{name} 속성이 같은지 비교하도록 정의했다.

따라서 \zsh{==}는 \zsh{True}가 된다.

정리하면 다음과 같다.

\begin{dhktable}
|연산자|비교 기준|
|:---|:---|
|is|두 표현식이 동일한 객체를 참조하는가?|
|==|객체가 정의한 동등성 비교 결과는 무엇인가?|
\end{dhktable}

특히 머신러닝에서는 Tensor 객체의 \zsh{==}가 원소별 비교를 수행한다는 점도 기억해 둘 필요가 있다.

따라서 Python 객체의 동일성 비교와 Tensor 값 비교는 구분해서 사용해야 한다.

\subsubsection{\zsh{type()}과 \zsh{isinstance()}의 차이}

이번에는 객체의 타입을 검사하는 두 방법을 비교해 보자.

\begin{dhknote}
\begin{dhksmry}[centered]
type() - 객체의 실제 타입 확인
\end{dhksmry}
\end{dhknote}

\begin{dhkcode}
class Person:
    pass

p = Person()
print(type(p))
\end{dhkcode}

출력:

\begin{dhkzsh}
<class '__main__.Person'>
\end{dhkzsh}

정확히 \zsh{Person} 클래스의 인스턴스인지 확인하려면 다음처럼 작성할 수 있다.

\begin{dhkcode}
print(type(p) is Person)
\end{dhkcode}

출력:

\begin{dhkzsh}
True
\end{dhkzsh}

\begin{dhknote}
\begin{dhksmry}[centered]
isinstance() - 클래스 관계를 고려한 검사
\end{dhksmry}
\end{dhknote}

\begin{dhkcode}
print(isinstance(p, Person))
\end{dhkcode}

출력:

\begin{dhkzsh}
True
\end{dhkzsh}

그런데 상속이 등장하면 차이가 발생한다.

\begin{dhkcode}
class Person:
    pass

class Student(Person):
    pass

s = Student()
\end{dhkcode}

여기서 \zsh{Student}는 \zsh{Person}을 상속받는다.

다음 결과를 확인해 보자.

\begin{dhkcode}
print(type(s) is Person)
print(type(s) is Student)

print(isinstance(s, Person))
print(isinstance(s, Student))
\end{dhkcode}

출력:

\begin{dhkzsh}
False
True
True
True
\end{dhkzsh}

왜 이런 결과가 발생할까?

\zsh{s}의 실제 타입은 \zsh{Student}다.

따라서 다음 비교는 \zsh{False}다.

\begin{dhkcode}
type(s) is Person
\end{dhkcode}

하지만 \zsh{Student}는 \zsh{Person}을 상속하므로 \zsh{s}는 \zsh{Person}의 인스턴스이기도 하다.

따라서 다음 비교는 \zsh{True}다.

\begin{dhkcode}
isinstance(s, Person)
\end{dhkcode}

정리하면 다음과 같다.

\begin{dhktable}
|표현식|	의미|
|:---|:---|
|type(s) is Person|	실제 타입이 정확히 Person인가?|
|type(s) is Student	|실제 타입이 정확히 Student인가?|
|isinstance(s, Person)	|Person 또는 그 하위 클래스의 인스턴스인가?|
|isinstance(s, Student)	|Student 또는 그 하위 클래스의 인스턴스인가?|
\end{dhktable}

\zsh{isinstance()}에는 추상 기반 클래스나 사용자 정의 인스턴스 검사 규칙 등 추가적인 동작도 있지만,

현재는 상속 관계를 고려한다는 점을 이해하면 충분하다.

이 개념은 이후 PyTorch에서 중요하게 사용된다.

\subsubsection{클래스 정의를 다시 실행하면 어떻게 되는가?}

이번에는 이름과 객체의 바인딩 관점에서 조금 더 깊게 살펴보자.

다음 코드를 실행한다고 가정하자.

\begin{dhkcode}
class Person:
    pass

p = Person()
\end{dhkcode}

이 시점의 관계는 다음과 같다.

\begin{dhkdiagram}[in=6]
node per (1,1) :: Person
node c (1,3) :: 클래스 객체 A
node t (3,3) :: type
node i (5,3) :: 인스턴스 p

edge per -> c
edge i -> t
edge t -> c
\end{dhkdiagram}

그런데 같은 이름으로 클래스를 다시 정의하면 어떻게 될까?

\begin{dhkcode}
class Person:
    pass
\end{dhkcode}

Python은 기존 클래스 객체를 수정하는 것이 아니라, 새로운 클래스 객체를 생성하고

\zsh{Person}이라는 이름을 그 객체에 다시 바인딩한다.

따라서 기존 인스턴스 \zsh{p}는 이전 클래스 객체를 타입으로 유지한다.

실제로 확인해 보자.

\begin{dhkcode}
class Person:
    pass

p = Person()

OldPerson = Person

class Person:
    pass


print(OldPerson is Person)
print(type(p) is OldPerson)
print(type(p) is Person)
\end{dhkcode}

출력:

\begin{dhkzsh}
False
True
False
\end{dhkzsh}

실행 과정을 나타내면 다음과 같다.

처음 클래스 정의 후

\begin{dhkdiagram}[in=6]
node per (1,1) :: Person
node c (1,3) :: 클래스 객체 A
edge per ->c
\end{dhkdiagram}

인스턴스 생성 후

\begin{dhkdiagram}[in=6]
node per (1,1) :: Person
node c (1,3) :: 클래스 객체 A
node i (3,3) :: 인스턴스 p
edge per ->c
edge i ->c
\end{dhkdiagram}

클래스 재정의 후

\begin{dhkdiagram}[in=6]
node op (1,1) :: OldPerson
node c (1,3) :: 클래스 객체 A
node i (3,3) :: 인스턴스 p
node per (5,1) :: Person
node b (5,3) :: 클래스 객체 B

edge op -> c
edge i -> c
edge per -> b
\end{dhkdiagram}

여기서 \zsh{Person}이라는 이름은 새로운 클래스 객체 B를 참조하지만,

기존 인스턴스 \zsh{p}의 타입은 여전히 클래스 객체 A다.

이 개념은 Jupyter Notebook에서 클래스를 여러 번 정의할 때 특히 중요하다.

예를 들어 Notebook에서 모델 클래스를 정의한 셀을다시 실행하면 새로운 클래스 객체가 만들어질 수 있다.

하지만 이전에 생성한 모델 인스턴스가 자동으로 새로운 클래스의 인스턴스로 바뀌는 것은 아니다.

따라서 다음과 같은 상황이 발생할 수 있다.

\begin{dhkcode}
model = MyModel()
\end{dhkcode}

이후 \zsh{MyModel} 클래스 정의 셀을 다시 실행하더라도 기존 \zsh{model}은

이전 클래스 정의를 기반으로 생성된 인스턴스다.

이런 경우에는 새 클래스 정의를 반영하려면 일반적으로 모델 인스턴스를 다시 생성해야 한다.

\subsubsection{PyTorch와의 연결 - 모델 클래스와 모델 인스턴스}

이제 PyTorch 코드에 적용해 보자.

\begin{dhkcode}
import torch.nn as nn

class MyModel(nn.Module):
    def __init__(self):
        super().__init__()

        self.linear = nn.Linear(10,5)

    def forward(self,x):
        return self.linear(x)
\end{dhkcode}

이 코드를 실행하면 \zsh{MyModel}이라는 클래스 객체가 생성된다.

아직 모델 인스턴스를 생성한 것은 아니다.

다음 코드를 실행하면

\begin{dhkcode}
model = MyModel()
\end{dhkcode}

\zsh{MyModel} 클래스 객체가 호출되고 모델 인스턴스가 생성된다.

여기서 중요한 차이가 있다.

\begin{dhknote}
\begin{dhksmry}[centered]
모델 클래스
\end{dhksmry}
\end{dhknote}

\begin{dhkcode}
MyModel
\end{dhkcode}

모델의 구조와 동작을 정의하는 클래스 객체다.

\begin{dhknote}
\begin{dhksmry}[centered]
모델 인스턴스
\end{dhksmry}
\end{dhknote}

\begin{dhkcode}
model = MyModel()
\end{dhkcode}

실제 \zsh{nn.Linear} 계층과 파라미터를 포함하는 모델 객체다.

\begin{dhknote}
\begin{dhksmry}[centered]
여러 모델 인스턴스 생성
\end{dhksmry}
\end{dhknote}

\begin{dhkcode}
model1 = MyModel()
model2 = MyModel()
\end{dhkcode}

두 모델은 같은 클래스의 인스턴스지만 서로 다른 객체다.

\begin{dhkcode}
print(model1 is model2)
\end{dhkcode}

출력:

\begin{dhkzsh}
False
\end{dhkzsh}

또한 각각 \zsh{nn.Linear} 계층도 별도로 생성된다.

\begin{dhkcode}
print(model1.linear is model2.linear)
\end{dhkcode}

출력:

\begin{dhkzsh}
False
\end{dhkzsh}

따라서 일반적으로 두 모델은 서로 독립적인 파라미터를 가진다.

\begin{dhkdiagram}[in=6]
node mm (1,2) :: MyModel
node c (3,2) :: 클래스 객체
node m1 (4,1) :: MyModel()
node m2 (4,3) :: MyModel()
node mm1 (6,1) :: model1
node mm2 (6,3) :: model2
node la (8,1) :: nn.Linear A
node lb (8,3) :: nn.Linear B
node pa (10,1) :: Parameters A
node pb (10,3) :: Parameters B
edge mm ->c
split c -> m1,m2
edge m1 -> mm1
edge mm1 -> la
edge la -> pa
edge m2 -> mm2
edge mm2 -> lb
edge lb -> pb
\end{dhkdiagram}

여기서 중요한 것은 클래스 객체는 공유하지만,

인스턴스 내부에 생성된 계층과 파라미터는 별도로 존재한다는 점이다.

이는 딥러닝 실험에서 중요한 의미를 가진다.

예를 들어 동일한 모델 아키텍처를 사용하면서 서로 다른 초기 파라미터를 가진 모델을 여러 개 생성할 수 있다.

\begin{dhkcode}
model1 = MyModel()
model2 = MyModel()
\end{dhkcode}

두 모델은 같은 아키텍처를 가지지만, 별도의 파라미터 객체를 가진다.

초기 가중치의 실제 값이 서로 다른지는 난수 상태와 초기화 방식에 따라 달라진다.

이러한 구조는 여러 모델을 비교하거나 서로 다른 학습 실험을 수행할 때 활용된다.

\subsubsection{핵심 정리}

이번 학습에서는 클래스 객체와 인스턴스 객체의 관계를 살펴보았다.

핵심 내용을 정리하면 다음과 같다.

\begin{dhknote}
\begin{dhksmry}[centered]
클래스 객체와 인스턴스 객체는 서로 다른 객체다.
\end{dhksmry}
\end{dhknote}

\begin{dhkcode}
Person
Person()
\end{dhkcode}

\zsh{Person}은 클래스 객체를 참조하고, \zsh{Person()}은 클래스 객체를 호출하는 표현식이다.

\begin{dhknote}
\begin{dhksmry}[centered]
인스턴스는 클래스 객체를 호출하여 생성할 수 있다.
\end{dhksmry}
\end{dhknote}

\begin{dhkcode}
p = Person()
\end{dhkcode}

일반적인 클래스에서는 새로운 인스턴스가 생성되고 \zsh{p}에 바인딩된다.

\begin{dhknote}
\begin{dhksmry}[centered]
하나의 클래스에서 여러 인스턴스를 생성할 수 있다.
\end{dhksmry}
\end{dhknote}

\begin{dhkcode}
p1 = Person()
p2 = Person()
\end{dhkcode}

일반적으로 두 인스턴스는 서로 다른 객체다.

\begin{dhknote}
\begin{dhksmry}[centered]
인스턴스를 생성한다고 클래스 객체가 복사되는 것은 아니다.
\end{dhksmry}
\end{dhknote}

\begin{dhkcode}
type(p1) is type(p2)
\end{dhkcode}

같은 클래스에서 생성한 인스턴스들은 동일한 클래스 객체를 타입으로 가진다.

\begin{dhknote}
\begin{dhksmry}[centered]
이름의 재바인딩은 객체의 재생성과 다르다.
\end{dhksmry}
\end{dhknote}

\begin{dhkcode}
p2 = p1
\end{dhkcode}

이 코드는 새로운 인스턴스를 생성하지 않는다.

\begin{dhknote}
\begin{dhksmry}[centered]
is와 ==는 서로 다른 비교 연산이다.
\end{dhksmry}
\end{dhknote}

\begin{dhkcode}
p1 is p2
p1 == p2
\end{dhkcode}

\zsh{is}는 객체의 동일성을 검사하고, \zsh{==}는 객체의 동등성을 검사한다.

\begin{dhknote}
\begin{dhksmry}[centered]
type()과 isinstance()는 검사 목적이 다르다.
\end{dhksmry}
\end{dhknote}

\begin{dhkcode}
type(p) is Person
isinstance(p, Person)
\end{dhkcode}

전자는 실제 타입을 정확히 비교하고, 후자는 일반적으로 상속 관계까지 고려한다.

\begin{dhknote}
\begin{dhksmry}[centered]
클래스를 다시 정의하면 새로운 클래스 객체가 생성된다.
\end{dhksmry}
\end{dhknote}

기존 인스턴스는 이전 클래스 객체를 자신의 타입으로 유지한다.

\begin{dhknote}
\begin{dhksmry}[centered]
PyTorch 모델도 클래스와 인스턴스를 구분해야 한다.
\end{dhksmry}
\end{dhknote}

\begin{dhkcode}
MyModel
MyModel()
\end{dhkcode}

\zsh{MyModel}은 모델 클래스 객체이고,

\zsh{MyModel()}은 일반적으로 실제 계층과 파라미터를 가진 모델 인스턴스를 생성한다.

\begin{dhknote}
\begin{dhksmry}[centered]
2.3에서 반드시 이해해야 할 결론
\end{dhksmry}
\end{dhknote}

이번 학습의 핵심 구조는 다음과 같다.

\begin{dhkdiagram}[in=12]
node c (1,1) :: Class Object
node i (3,1) :: Instance Object
node n (5,1) :: Name Binding
edge c->i :: 호출
edge i ->n :: 반환
\end{dhkdiagram}

그리고 하나의 클래스 객체를 기반으로 여러 인스턴스가 생성될 수 있다.

\begin{dhkdiagram}[in=10]
node per (3,1) :: Person
node p1 (1,3) :: p1
node p2 (3,3) :: p2
node p3 (5,3) :: p3

edge per -> p1
edge per -> p2
edge per -> p3
\end{dhkdiagram}

여기서 \zsh{Person}은 클래스 객체이고,

\zsh{p1, p2, p3}는 각각 해당 클래스의 인스턴스를 참조하는 이름이다.

클래스는 인스턴스의 타입과 동작을 정의하고, 인스턴스는 그 클래스를 기반으로 생성된 개별 객체다.

다음 \zsh{2.4. 인스턴스 속성과 네임스페이스}에서는

지금까지 구분한 클래스 객체와 인스턴스 객체를 기반으로

\zsh{self.name}이 실제로 어디에 저장되는지,

\zsh{__dict__}가 어떤 역할을 하는지,

Python이 속성을 어떻게 찾아가는지 살펴볼 것이다.